% !TEX root = ../main.tex
\section{Introduzione}
\label{sec:introduzione}

\subsection{Rilevanza del problema e obiettivo}
\subsubsection{Perché fidarsi di una dipendenza può esporre a un attacco}

Lo sviluppo software moderno si basa spesso su componenti realizzati da terzi,
così da riutilizzarne le funzionalità. Questa pratica introduce un confine di
fiducia: una dipendenza può diventare un vettore d'attacco se chi la mantiene o
il processo con cui viene distribuita viene compromesso.

\subsubsection{Casi selezionati e domanda guida}

La relazione confronta due episodi: il phishing che nel settembre 2025 portò
alla compromissione dell'account npm di un manutentore di \texttt{chalk} e
\texttt{debug}, e l'ondata di adozioni e aggiornamenti malevoli nell'Arch User
Repository (AUR) segnalata nel giugno 2026~\cite{aikido2025,arch-aur-incident-2026}.
Il primo caso sfruttò l'autorità di pubblicazione di un account; il secondo la
fiducia nelle ricette e negli script di installazione dei pacchetti AUR.

L'obiettivo è confrontare i passaggi di fiducia abusati, i rischi che ne
derivano e le misure che possono ridurli. La domanda guida è: \textbf{in che
modo la compromissione di un account di pubblicazione o la modifica di una
ricetta d'installazione può trasformare una dipendenza legittima in un vettore
d'attacco?}
